Define the "span" $s(X)$ of a subset $X$ of $K$ as the algebraic closure of $k(X)$ in $K$ (the algebraic closure of $k$ in $K$, if $X$ is empty). Then the following conditions are satisfied:

\begin{enumerate}
\item[$(S_1)$] If $X \subset Y$, then $s(X) \subset s(Y)$.
\item[$(S_2)$] If $x \in s(X)$, then there exists a finite subset $Y$ of $X$ such that $x \in s(Y)$.
\item[$(S_3)$] $X \subset s(X)$ for all subsets $X$ of $K$.
\item[$(S_4)$] $s(s(X)) = s(X)$ (this simply expresses the transitivity of algebraic dependence).
\item[$(S_5)$] The relations $y \in s(X, x)$ and $y \notin s(X)$ imply $x \in s(X, y)$.
\end{enumerate}
